\textbf{Proof.}
Fix \(P\in\mathfrak P_r(N_0,K_T,K_D)\), and set
\(q=\min\{\beta+1,2\}\).  By the hazard identity and endpoint
normalization,
\[
 \theta(t)=S^0_{1,\star}(t)
 =\exp\left\{-\int_0^t\lambda^0_{1,\star}(u)\,du\right\}.
\]
The bound \(\|\lambda^0_{1,\star}\|_{C^\beta(\mathcal I)}\le L\), the
fixed horizon, and the product H\"older inequality applied to
\(\theta'=-\lambda^0_{1,\star}\theta\) imply
\[
 \|\theta\|_{C^{\beta+1}(\mathcal I)}
 \le C(L,\bar\tau,\beta).                                   \tag{1}
\]
Let \(I_{\mathcal G}\theta\) be the piecewise-linear interpolant of the
grid values of \(\theta\).  On a grid interval \([a,b]\) of length at most
\(h=h_{N_0,K_T,K_D}\), the first-order Taylor formula with H\"older
remainder when \(q<2\), and the ordinary second-order interpolation
remainder when \(q=2\), give
\[
 \sup_{t\in[a,b]}|I_{\mathcal G}\theta(t)-\theta(t)|
 \le C(L,\bar\tau,\beta)|b-a|^q.
\]
Taking the maximum over intervals yields
\[
 \|I_{\mathcal G}\theta-\theta\|_\infty
 \le c_{\mathrm{int}}h^q.                                  \tag{2}
\]

By construction, \(\widehat\theta_r\) is linear between adjacent grid
points.  Hence \(\widehat\theta_r-I_{\mathcal G}\theta\) is also linear
there, and the maximum of its absolute value is attained at an endpoint.
Consequently,
\[
 \|\widehat\theta_r-\theta\|_{\infty,\mathcal I}
 \le \max_{g\in\mathcal G}|\widehat\theta_r(g)-\theta(g)|
 +\|I_{\mathcal G}\theta-\theta\|_\infty.                  \tag{3}
\]
The reverse inequality between the grid maximum and the continuous
supremum follows from \(\mathcal G\subset\mathcal I\).  Combining (2)--(3)
proves the asserted two-sided interpolation inequality.  Finally, the mesh
condition in \(\mathrm{def:inference\mbox{-}grid}\) is \(h^q\le\rho_r\),
so the allowance is \(O(\rho_r)\). \(\square\)